\honest{The Scales of Justice, the Notebook of Rationality}{Eliezer Yudkowsky}{2007}

Lady Justice carries scales, and on scales whatever pulls one side down pushes the other up. That is convenient, and ``usually a gross distortion.'' People treat debate like a sport: every point scored against one side counts for the other, and the side with more points is right about everything.

A study called ``The Affect Heuristic in Judgments of Risks and Benefits'' asked whether people merge a technology's risks and benefits into one good or bad feeling. Suppose a reactor makes less waste than other designs but is more likely to melt down. One fact is a point for the reactor and one a point against, but they are not opposed. ``The facts don't know whose side they're on.'' A reactor that is passively safe need not also make less waste or cheaper power; these are all good, but not the same good. How much waste a design makes and how likely it is to melt down are two physical questions with two answers.

People judge by overall feeling, and ``If you tell people a reactor design produces less waste, they rate its probability of meltdown as lower.'' That gets physical questions wrong by treating facts as soldiers, any one of which can be sent against any soldier on the other side. \nb{The study did not test waste against meltdown. It gave people a paragraph about the risks or the benefits of nuclear power in general, and their ratings of the attribute not mentioned generally moved the other way. The study's authors add that in the world risks and benefits tend to go together, which is why the pattern counts as an error.}

Scales do suit one kind of question: a single factual question with two possible answers, such as whether John Smith killed John Doe. E. T. Jaynes taught that evidence only moves probability between hypotheses. \nb{Jaynes wrote that, unless the facts are impossible on a hypothesis, ``it is meaningless to ask how much those facts tend `in themselves' to confirm or refute'' it.} For any more complex issue, Justice ``should relinquish her scales or relinquish her sword.'' \nb{A decision still weighs facts against each other. In ``Policy Debates Should Not Appear One-Sided,'' ten days earlier, the author's example of real tough-mindedness was keeping the shops open ``because we did this cost-benefit calculation.''}

Lady Rationality carries a notebook instead, and writes down all the facts that are on no one's side.
